%
% this file is encoded in utf-8
% v1.7
% do not change the content of this file
% unless the thesis layout rule is changed
% 無須修改本檔內容，除非校方修改了
% 封面、書名頁、中文摘要、英文摘要、誌謝、目錄、表目錄、圖目錄、符號說明
% 等頁之格式
% this file is encoded in utf-8
%v1.7

% make the line spacing in effect
\renewcommand{\baselinestretch}{\mybaselinestretch}
\large % it needs a font size changing command to be effective

% default variables definitions
% 注意！！此處只是預設值，不需更改此處
% 請更改 my_names.tex 內容
\newcommand\cTitle{論文題目}
\newcommand\eTitle{MY THESIS TITLE}
\newcommand\myCname{王鐵雄}
\newcommand\myEname{Aron Wang}
\newcommand\myStudentID{M9315048}
\newcommand\advisorCnameA{南宮明博士}
\newcommand\advisorEnameA{Dr.~Ming Nangong}
\newcommand\advisorCnameB{李斯坦博士}
\newcommand\advisorEnameB{Dr.~Stein Lee}
\newcommand\advisorCnameC{徐　石博士}
\newcommand\advisorEnameC{Dr.~Sean~Hsu}
\newcommand\univCname{國立台灣科技大學}
\newcommand\univEname{National Taiwan University of science and technology}
\newcommand\deptCname{光電工程研究所}
\newcommand\fulldeptEname{Graduate School of Electro-Optical Engineering}
\newcommand\deptEname{Electro-optical Engineering}
\newcommand\collEname{College of Engineering}
\newcommand\degreeCname{碩士}
\newcommand\degreeEname{Master of Science}
\newcommand\cYear{九十四}
\newcommand\cMonth{六}
\newcommand\cDay{十}
%\newcounter{eYear}
\newcommand\eYear{2006}
\newcommand\eMonth{June}
\newcommand\ePlace{Chungli, Taoyuan, Taiwan}


 % user's names; to replace those default variable definitions
%
% this file is encoded in utf-8
% v1.7
% 填入你的論文題目、姓名等資料
% 如果題目內有必須以數學模式表示的符號，請用 \mbox{} 包住數學模式，如下範例
% 如果中文名字是單名，與姓氏之間建議以全形空白填入，如下範例
% 英文名字中的稱謂，如 Prof. 以及 Dr.，其句點之後請以不斷行空白~代替一般空白，如下範例
% 如果你的指導教授沒有如預設的三位這麼多，則請把相對應的多餘教授的中文、英文名
%    的定義以空的大括號表示
%    如，\renewcommand\advisorCnameB{}
%          \renewcommand\advisorEnameB{}
%          \renewcommand\advisorCnameC{}
%          \renewcommand\advisorEnameC{}

% 論文題目 (中文)
\renewcommand\cTitle{%我的碩士論文題目 
校舍耐震資料庫之資料探勘
}

% 論文題目 (英文)
\renewcommand\eTitle{%My Thesis Title  
Data Mining on Aseismic School Building Database 
%My Thesis Title  \mbox{$\cal{H}_\infty$} and \mbox{Al$_x$Ga$_{1-x}$As}
}

% 我的姓名 (中文)
\renewcommand\myCname{高偉格}

% 我的姓名 (英文)
\renewcommand\myEname{Kao, Wei-Ko}

%我的學號
\renewcommand\myStudentID{D9505501}

% 指導教授A的姓名 (中文)
\renewcommand\advisorCnameA{陳鴻銘 博士}

% 指導教授A的姓名 (英文)
\renewcommand\advisorEnameA{Dr. Hung-Ming Cheng}

% 指導教授B的姓名 (中文)
\renewcommand\advisorCnameB{}

% 指導教授B的姓名 (英文)
\renewcommand\advisorEnameB{}

% 指導教授C的姓名 (中文)
\renewcommand\advisorCnameC{}

% 指導教授C的姓名 (英文)
\renewcommand\advisorEnameC{}

% 校名 (中文)
\renewcommand\univCname{國立台灣科技大學}

% 校名 (英文)
%\renewcommand\univEname{National Taiwan University of science and technology}

% 系所名 (中文)
\renewcommand\deptCname{營~建~工~程~系}

% 系所全名 (英文)
%\renewcommand\fulldeptEname{Graduate School of Electro-Optical Engineering}

% 系所短名 (英文, 用於書名頁學位名領域)
%\renewcommand\deptEname{Electro-Optical Engineering}

% 學院英文名 (如無，則以空的大括號表示)
%\renewcommand\collEname{College of Electrical and Communication Engineering}

% 學位名 (中文)
\renewcommand\degreeCname{博士學位}

% 學位名 (英文)
%\renewcommand\degreeEname{Master of Science}

% 口試年份 (中文、民國)
\renewcommand\cYear{一零二}

% 口試月份 (中文)
\renewcommand\cMonth{七} 

% 口試月份 (中文)
\renewcommand\cDay{七} 

% 口試年份 (阿拉伯數字、西元)
%\renewcommand\eYear{2009} 

% 口試月份 (英文)
%\renewcommand\eMonth{July}

% 學校所在地 (英文)
%\renewcommand\ePlace{Taipei, Taiwan}

%畢業級別；用於書背列印；若無此需要可忽略
\newcommand\GraduationClass{95}

%%%%%%%%%%%%%%%%%%%%%%
\newcommand\itsempty{}
%%%%%%%%%%%%%%%%%%%%%%%%%%%%%%%
%       ntust cover 封面
%%%%%%%%%%%%%%%%%%%%%%%%%%%%%%%
%
\begin{titlepage}
% no page number
% next page will be page 1

% aligned to the center of the page
\begin{center}
% font size (relative to 12 pt):
% \large (14pt) < \Large (18pt) < \LARGE (20pt) < \huge (24pt)< \Huge (24 pt)
%



\begin{figure}[htbp]
	\begin{minipage}[t]{5cm} 
		\raggedright
		\includegraphics[width=1.1in]{frontpages/ntust_logo.eps}
		\label{fig:ntust_logo}
		\vspace{0.5cm}
	\end{minipage}% 
	\begin{minipage}[b]{0.5\textwidth} 
	\centering
	\makebox[3cm][c]{\Huge{\univCname}}\\  %顯示中文校名
	\vspace{0.5cm}
	\makebox[3cm][c]{\Huge{\deptCname}}\\ %顯示中文系所名
	\vspace{0.5cm}
	\makebox[3cm][c]{\Huge{\degreeCname 論文}}\\ %顯示論文種類 (中文)
	\vspace{0.5cm}
	\end{minipage}%
\\ 
\rule{16cm}{3pt}
\end{figure}
%\hfill


%\makebox[8cm][s]{\textbf{\Huge{\degreeCname 論文（初稿）}}}\\ %顯示論文種類 (中文)
\vspace{1cm}
%
% Set the line spacing to single for the titles (to compress the lines)
\renewcommand{\baselinestretch}{1}   %行距 1 倍
%\large % it needs a font size changing command to be effective
\textbf{\LARGE{\cTitle}}\\  % 中文題目
%
\vspace{1cm}
%
\Large{\eTitle}\\ %英文題目
\vspace{5cm}
% \makebox is a text box with specified width;
% option s: stretch
% use \makebox to make sure
% 「研究生：」 與「指導教授：」occupy the same width
\hspace{4.5cm} \makebox[3cm][s]{\Large{研 究 生：}}
\Large{\myCname}  % 顯示作者中文名
\hfill \makebox[1cm][s]{}\\
%
\vspace{0.3cm}
\hspace{4.5cm} \makebox[3cm][s]{\Large{學號：}}
\Large{\myStudentID}  %顯示指導教授A中文名
\hfill \makebox[1cm][s]{}\\
%
\vspace{1cm}
\hspace{4.5cm} \makebox[3cm][s]{\Large{指導教授：}}
\Large{\advisorCnameA}  %顯示指導教授A中文名
\hfill \makebox[1cm][s]{}\\
%
% 判斷是否有共同指導的教授 B
\ifx \advisorCnameB  \itsempty
\relax % 沒有 B 教授，所以不佔版面，不印任何空白
\else
% 共同指導的教授 B
\hspace{4.5cm} \makebox[3cm][s]{}
\Large{\advisorCnameB}  %顯示指導教授B中文名
\hfill \makebox[1cm][s]{}\\
\fi
%
% 判斷是否有共同指導的教授 C
\ifx \advisorCnameC  \itsempty
\relax % 沒有 C 教授，所以不佔版面，不印任何空白
\else
% 共同指導的教授 C
\hspace{4.5cm} \makebox[3cm][s]{}
\Large{\advisorCnameC}  %顯示指導教授C中文名
\hfill \makebox[1cm][s]{}\\
\fi
%
\vfill
\makebox[10cm][s]{\Large{中華民國\cYear 年\cMonth 月\cDay 日}}%
%
\end{center}
% Resume the line spacing to the desired setting
\renewcommand{\baselinestretch}{\mybaselinestretch}   %恢復原設定
% it needs a font size changing command to be effective
% restore the font size to normal
\normalsize
\end{titlepage}
%%%%%%%%%%%%%%

%% 從摘要到本文之前的部份以小寫羅馬數字印頁碼
% 但是從「書名頁」(但不印頁碼) 就開始計算
%\setcounter{page}{1}
\pagenumbering{Roman}
%\pagenumbering{arabic}
%%%%%%%%%%%%%%%%%%%%%%%%%%%%%%%
%       指導教授推薦書 
%%%%%%%%%%%%%%%%%%%%%%%%%%%%%%%
%
% insert the printed standard form when the thesis is ready to bind
% 在口試完成後，再將已簽名的推薦書放入以便裝訂
% create an entry in table of contents for 推薦書
% 目前送出空白頁
%\newpage{\thispagestyle{empty}\addcontentsline{toc}{chapter}{\nameInnerCover}\mbox{}\clearpage}%
%\newpage

% 判斷是否要浮水印？
\ifx\mywatermark\undefined 
  \thispagestyle{empty}  % 無頁碼、無 header (無浮水印)
\else
  \thispagestyle{EmptyWaterMarkPage} % 無頁碼、有浮水印
\fi

%%%%%%%%%%%%%%%%%%%%%%%%%%%%%%%%%%%%%%%%%%%%%%%%%%%%%%%%%%%%%%%
%%no page number
%% create an entry in table of contents for 書名頁
%\addcontentsline{toc}{chapter}{\nameInnerCover}
%
%
%% aligned to the center of the page
%\begin{center}
%% font size (relative to 12 pt):
%% \large (14pt) < \Large (18pt) < \LARGE (20pt) < \huge (24pt)< \Huge (24 pt)
%% Set the line spacing to single for the titles (to compress the lines)
%\renewcommand{\baselinestretch}{1}   %行距 1 倍
%% it needs a font size changing command to be effective
%%中文題目
%\Large{\cTitle}\\ %%%%%
%\vspace{1cm}
%% 英文題目
%\Large{\eTitle}\\ %%%%%
%%\vspace{1cm}
%\vfill
%% \makebox is a text box with specified width;
%% option s: stretch
%% use \makebox to make sure
%% 「研究生：」 與「指導教授：」occupy the same width
%\makebox[3cm][s]{\large{研 究 生：}}
%\makebox[3cm][l]{\large{\myCname}} %%%%%
%\hfill
%\makebox[2cm][s]{\large{Student: }}
%\makebox[5cm][l]{\large{\myEname}}\\ %%%%%
%%
%%\vspace{1cm}
%%
%\makebox[3cm][s]{\large{指導教授：}}
%\makebox[3cm][l]{\large{\advisorCnameA}} %%%%%
%\hfill
%\makebox[2cm][s]{\large{Advisor: }}
%\makebox[5cm][l]{\large{\advisorEnameA}}\\ %%%%%
%%
%% 判斷是否有共同指導的教授 B
%\ifx \advisorCnameB  \itsempty
%\relax % 沒有 B 教授，所以不佔版面，不印任何空白
%\else
%%共同指導的教授B
%\makebox[3cm][s]{}
%\makebox[3cm][l]{\large{\advisorCnameB}} %%%%%
%\hfill
%\makebox[2cm][s]{}
%\makebox[5cm][l]{\large{\advisorEnameB}}\\ %%%%%
%\fi
%%
%% 判斷是否有共同指導的教授 C
%\ifx \advisorCnameC  \itsempty
%\relax % 沒有 C 教授，所以不佔版面，不印任何空白
%\else
%%共同指導的教授C
%\makebox[3cm][s]{}
%\makebox[3cm][l]{\large{\advisorCnameC}} %%%%%
%\hfill
%\makebox[2cm][s]{}
%\makebox[5cm][l]{\large{\advisorEnameC}}\\ %%%%%
%\fi
%%
%% Resume the line spacing to the desired setting
%\renewcommand{\baselinestretch}{\mybaselinestretch}   %恢復原設定
%\large %it needs a font size changing command to be effective
%%
%\vfill
%\makebox[4cm][s]{\large{\univCname}}\\% 校名
%\makebox[6cm][s]{\large{\deptCname}}\\% 系所名
%\makebox[3cm][s]{\large{\degreeCname 論文}}\\% 學位名
%%
%%\vspace{1cm}
%\vfill
%\large{A Thesis}\\%
%\large{Submitted to }%
%%
%\large{\fulldeptEname}\\%系所全名 (英文)
%%
%%
%\ifx \collEname  \itsempty
%\relax % 沒有學院名 (英文)，所以不佔版面，不印任何空白
%\else
%% 有學院名 (英文)
%\large{\collEname}\\% 學院名 (英文)
%\fi
%%
%\large{\univEname}\\%校名 (英文)
%%
%\large{in Partial Fulfillment of the Requirements}\\
%%
%\large{for the Degree of}\\
%%
%\large{\degreeEname}\\%學位名(英文)
%
%\large{in}\\
%%
%\large{\deptEname}\\%系所短名(英文；表明學位領域)
%%
%\large{\eMonth\ \eYear}\\%月、年 (英文)
%%
%\large{\ePlace}% 學校所在地 (英文)
%\vfill
%\large{中華民國}%
%\large{\cYear}% %%%%%
%\large{年}%
%\large{\cMonth}% %%%%%
%\large{月}\\
%\end{center}
%% restore the font size to normal
%\normalsize
%\clearpage


%%%%%%%%%%%%%%%%%%%%%%%%%%%%%%%%%%%%%%%%%%%%%%%%%%%%%%%%%%%%%%%%%%%%%
%%%%%%%%%%%%%%%%%%%%%%%%%%%%%%%
%       論文口試委員審定書 (計頁碼，但不印頁碼) 
%%%%%%%%%%%%%%%%%%%%%%%%%%%%%%%
%
% insert the printed standard form when the thesis is ready to bind
% 在口試完成後，再將已簽名的審定書放入以便裝訂
% create an entry in table of contents for 審定書
% 目前送出空白頁

%\newpage{\thispagestyle{empty}\addcontentsline{toc}{chapter}{\nameCommitteeForm}\mbox{}\clearpage}%


%%%%%%%%%%%%%%%%%%%%%%%%%%%%%%%
%       中文摘要 
%%%%%%%%%%%%%%%%%%%%%%%%%%%%%%%
%
\newpage
\thispagestyle{plain}  % 無 header，但在浮水印模式下會有浮水印
% create an entry in table of contents for 中文摘要
\addcontentsline{toc}{chapter}{\nameCabstract}

% aligned to the center of the page
\begin{center}
% font size (relative to 12 pt):
% \large (14pt) < \Large (18pt) < \LARGE (20pt) < \huge (24pt)< \Huge (24 pt)
% Set the line spacing to single for the names (to compress the lines)
\renewcommand{\baselinestretch}{1}   %行距 1 倍
% it needs a font size changing command to be effective
%\large{\cTitle}\\  %中文題目
%\vspace{0.83cm}
% \makebox is a text box with specified width;
% option s: stretch
% use \makebox to make sure
% each text field occupies the same width
%\makebox[1.5cm][s]{\large{學生：}}
%\makebox[3cm][l]{\large{\myCname}} %學生中文姓名
%\hfill
%
%\makebox[3cm][s]{\large{指導教授：}}
%\makebox[3cm][l]{\large{\advisorCnameA}} \\ %教授A中文姓名
%
% 判斷是否有共同指導的教授 B
\ifx \advisorCnameB  \itsempty
\relax % 沒有 B 教授，所以不佔版面，不印任何空白
\else
%共同指導的教授B
\makebox[1.5cm][s]{}
\makebox[3cm][l]{} %%%%%
\hfill
\makebox[3cm][s]{}
\makebox[3cm][l]{\large{\advisorCnameB}}\\ %教授B中文姓名
\fi
%
% 判斷是否有共同指導的教授 C
\ifx \advisorCnameC  \itsempty
\relax % 沒有 C 教授，所以不佔版面，不印任何空白
\else
%共同指導的教授C
\makebox[1.5cm][s]{}
\makebox[3cm][l]{} %%%%%
\hfill
\makebox[3cm][s]{}
\makebox[3cm][l]{\large{\advisorCnameC}}\\ %教授C中文姓名
\fi
%
%\vspace{0.42cm} %2013/06/06
\vspace{0.38cm}
%
%\large{\univCname}\large{\deptCname}\\ %校名系所名
%\vspace{0.83cm}
%\vfill
\makebox[2.5cm][c]{\Large{中文摘要}}\\
\end{center}
% Resume the line spacing to the desired setting
\renewcommand{\baselinestretch}{\mybaselinestretch}   %恢復原設定
%it needs a font size changing command to be effective
% restore the font size to normal
\normalsize
%%%%%%%%%%%%%

台灣國家實驗研究院地震工程研究中心（NCREE）在九二一大地震後，與教育部合作執行加速國中小老舊校舍及相關設備補強整建計畫，評估全台灣的各級學校校舍之耐震能力，在此計畫執行的過程中產生了大量的評估與調查資料，因此 NCREE 便建立了一個校舍耐震能力資料庫來收集各種相關的資料，收集了包括校舍的各種設計參數、材料強度、校舍現況及年齡、技師的評估與補強建議方案、實際補強的金額與補強方法等。此一資料庫所收集的校舍資料數量龐大，除了當初設計的目的之外，應該還潛藏難以由人直接判斷取得的知識（knowledge）、模式（pattern），而資料探勘（Data Mining）就是用來分析這種數量龐大的資料，可以從中找出潛藏知識的相關技術的統稱，本研究之目的即為利用資料探勘技術來發掘潛藏於此校舍耐震資料庫中的知識，並從資料探勘的四種主要分析方法：回歸、分類、分群、關聯出發，分別探討各種方法在此資料庫中有何可能的分析方向，有哪些可能的潛藏知識，並進行分析，最後得到了三個有用且可靠度足夠的關係模型，分別為校舍資訊與耐震能力之關係模型、校校舍資訊與破壞構件之關係模型以及校舍資訊與補強經費之關係模型。
	

%%%%%%%%%%%%%%%%%%%%%%%%%%%%%%%
%       英文摘要 
%%%%%%%%%%%%%%%%%%%%%%%%%%%%%%%
%
\newpage
\thispagestyle{plain}  % 無 header，但在浮水印模式下會有浮水印

% create an entry in table of contents for 英文摘要
\addcontentsline{toc}{chapter}{\nameEabstract}

% aligned to the center of the page
\begin{center}
% font size:
% \large (14pt) < \Large (18pt) < \LARGE (20pt) < \huge (24pt)< \Huge (24 pt)
% Set the line spacing to single for the names (to compress the lines)
\renewcommand{\baselinestretch}{1}   %行距 1 倍
%\large % it needs a font size changing command to be effective
%\large{\eTitle}\\  %英文題目
%\vspace{0.83cm}
% \makebox is a text box with specified width;
% option s: stretch
% use \makebox to make sure
% each text field occupies the same width
%\makebox[2cm][s]{\large{Student: }}
%\makebox[5cm][l]{\large{\myEname}} %學生英文姓名
%\hfill
%
%\makebox[2cm][s]{\large{Advisor: }}
%\makebox[5cm][l]{\large{\advisorEnameA}} \\ %教授A英文姓名
%
% 判斷是否有共同指導的教授 B
\ifx \advisorCnameB  \itsempty
\relax % 沒有 B 教授，所以不佔版面，不印任何空白
\else
%共同指導的教授B
\makebox[2cm][s]{}
\makebox[5cm][l]{} %%%%%
\hfill
\makebox[2cm][s]{}
\makebox[5cm][l]{\large{\advisorEnameB}}\\ %教授B英文姓名
\fi
%
% 判斷是否有共同指導的教授 C
\ifx \advisorCnameC  \itsempty
\relax % 沒有 C 教授，所以不佔版面，不印任何空白
\else
%共同指導的教授C
\makebox[2cm][s]{}
\makebox[5cm][l]{} %%%%%
\hfill
\makebox[2cm][s]{}
\makebox[5cm][l]{\large{\advisorEnameC}}\\ %教授C英文姓名
\fi
%
%\vspace{0.42cm}
%\large{Submitted to }\large{\fulldeptEname}\\  %英文系所全名
%
\ifx \collEname  \itsempty
\relax % 如果沒有學院名 (英文)，則不佔版面，不印任何空白
\else
% 有學院名 (英文)
%\large{\collEname}\\% 學院名 (英文)
\fi
%
%\large{\univEname}\\  %英文校名
%\vspace{0.83cm}
%\vfill
%
\large{ABSTRACT}\\
%\vspace{0.5cm}
\end{center}
% Resume the line spacing the desired setting
\renewcommand{\baselinestretch}{\mybaselinestretch}   %恢復原設定
%\large %it needs a font size changing command to be effective
% restore the font size to normal
\normalsize
%%%%%%%%%%%%%

台灣國家實驗研究院地震工程研究中心（NCREE）在九二一大地震後，與教育部合作評估全台灣的各級學校校舍之耐震能力，在此計畫執行的過程中產生了大量的評估與調查資料，因此 NCREE 便建立了一個校舍耐震能力資料庫來收集各種相關的資料，收集了包括校舍的各種設計參數、材料強度、校舍現況及年齡、技師的評估與補強建議方案、實際補強的金額與補強方法等。收集的校舍資料數量龐大，除了當初設計的目的之外，應該還潛藏難以由人直接判斷取得的知識（knowledge）、模式（pattern）。資料探勘（Data Mining）就是用來分析這種數量龐大的資料，從中找出潛藏的知識的相關技術的統稱，本研究之目的即為利用資料探勘技術來發掘潛藏於此校舍耐震資料庫中的知識，本研究從資料探勘的四種主要分析方法：回歸、分類、分群、關聯出發，分別探討各種方法在此資料庫中有何可能的分析方向，有哪些可能的潛藏知識，並進行分析，最後得到了三個有用的預測模型，分別為校舍耐震能力預測模型、校舍破壞模式預測模型以及校舍補強經費預估模型。
	
In general, the aseismic ability of buildings is analyzed using nonlinear models. To obtain buildings’ aseismic abilities, numerical models are constructed based on the structural configuration and material properties of buildings, and their stress responses and behaviors are simulated. This method is complex, time-consuming, and should only be conducted by professionals. In the past, soft computing techniques have been applied in the construction field to predict the particular stress responses and behaviors; however, only a few studies have been made to predict specific properties of entire buildings. In this study, a Weighted Genetic Programming system is developed to construct the relation models between the aseismic capacity of school buildings, and their basic design parameters. This is based on information from the database of school buildings, as well as information regarding the aseismic capacity of school buildings analyzed using complete nonlinear methods. This system can be further applied to predict the aseismic capacity of the school buildings.

%%%%%%%%%%%%%%%%%%%%%%%%%%%%%%%
%       誌謝 
%%%%%%%%%%%%%%%%%%%%%%%%%%%%%%%
%
% Acknowledgment
\newpage
\chapter*{\protect\makebox[5cm][s]{\nameAckn}} %\makebox{} is fragile; need protect

\addcontentsline{toc}{chapter}{\nameAckn}






%%%%%%%%%%%%%%%%%%%%%%%%%%%%%%%
%       目錄 
%%%%%%%%%%%%%%%%%%%%%%%%%%%%%%%
%
% Table of contents

\newpage


\renewcommand{\contentsname}{\protect\makebox[5cm][s]{\nameToc}}
%\makebox{} is fragile; need protect
\addcontentsline{toc}{chapter}{\nameToc}
\setcounter{tocdepth}{2}
\tableofcontents




%%%%%%%%%%%%%%%%%%%%%%%%%%%%%%%
%       圖目錄 
%%%%%%%%%%%%%%%%%%%%%%%%%%%%%%%
%
% List of Figures
\newcommand{\loflabel}{圖}
\newpage
\renewcommand{\numberline}[1]{\loflabel~#1\hspace*{1em}}
\renewcommand{\listfigurename}{\protect\makebox[5cm][s]{\nameTof}}
%\makebox{} is fragile; need protect
\addcontentsline{toc}{chapter}{\nameTof}
\listoffigures

%%%%%%%%%%%%%%%%%%%%%%%%%%%%%%%
%       表目錄 
%%%%%%%%%%%%%%%%%%%%%%%%%%%%%%%
%
% List of Tables
\newcommand{\lotlabel}{表}
\newpage
\renewcommand{\numberline}[1]{\lotlabel~#1\hspace*{1em}}
\renewcommand{\listtablename}{\protect\makebox[5cm][s]{\nameLot}}
%\makebox{} is fragile; need protect
\addcontentsline{toc}{chapter}{\nameLot}
\listoftables

%%%%%%%%%%%%%%%%%%%%%%%%%%%%%%%
%       演算法目錄 
%%%%%%%%%%%%%%%%%%%%%%%%%%%%%%%
%
% List of Figures
%\newpage
%\renewcommand{\listalgorithmname}{\protect\makebox[5cm][s]{\nameToa}}
%\makebox{} is fragile; need protect
%\addcontentsline{toc}{chapter}{\nameToa}
%\listofalgorithms


%%%%%%%%%%%%%%%%%%%%%%%%%%%%%%%
%       符號說明 
%%%%%%%%%%%%%%%%%%%%%%%%%%%%%%%
%
% Symbol list
% define new environment, based on standard description environment
% adapted from p.60~64, <<The LaTeX Companion>>, 1994, ISBN 0-201-54199-8
%\newcommand{\SymEntryLabel}[1]%
% {\makebox[3cm][l]{#1}}
%
%\newenvironment{SymEntry}
%   {\begin{list}{}%
%       {\renewcommand{\makelabel}{\SymEntryLabel}%
%        \setlength{\labelwidth}{3cm}%
%        \setlength{\leftmargin}{\labelwidth}%
%        }%
%   }%
%   {\end{list}}
%%
%\newpage
%\chapter*{\protect\makebox[5cm][s]{\nameSlist}} %\makebox{} is fragile; need protect
%\addcontentsline{toc}{chapter}{\nameSlist}
%%
% this file is encoded in utf-8
% v1.7
%  各符號以 \item[] 包住，然後接著寫說明
% 如果符號是數學符號，應以數學模式表示，以取得正確的字體
% 如果符號本身帶有方括號，則此符號可以用大括號 {} 包住保護
\begin{SymEntry}

\item[$Af$]
總樓地板面積

\item[$A_C$]
校舍破壞地表加速度

\item[$A_D$]
校舍所需承受之~475~年週期最大地震之地表加速度

\item[$CDR$]
耐震需求比

\item[$C$]
訓練模型時之懲罰值（Cost）參數

\item[$\bar{C}$]
WGP~決策樹中之常數端點

\item[$C_E$]
$Is$~值轉換之~$A_C$~等意義數值

\item[$ClaAc$]
一樓教室柱總斷面積

\item[$CorAc$]
一樓走廊外柱總斷面積

\item[$D\_isR$]
校舍耐震資料庫中紀錄校舍是否需要補強之屬性

\item[$E$]
初步評估之基本耐震性能

\item[$F$]
WGP~運算樹中，代表節點之運算子方程式

\item[$f_i$]
WGP~運算樹中，第~$i$~個節點之運算子方程式

\item[$I$]
用途係數

\item[$InsAc$]
一樓隔間柱總斷面積

\item[$Is$]
初步評估之耐震指標

\item[$K$]
K-means~分群方法之初始群集數

\item[$N$]
資料總數

\item[$NF$]
校舍樓層數

\item[$NI$]
輸入參數的數量

\item[$NL$]
運算層的層數

\item[$N_g$]
需要最佳化的基因數量

\item[$O$]
偏移變數

\item[$P_i$]
WGP~運算樹中之第~$i$~個輸入屬性

\item[$P_1$]
樓層數，WGP~建立之校舍耐震能力關係模型參數之一

\item[$P_2$]
鋼筋強度，WGP~建立之校舍耐震能力關係模型參數之一

\item[$P_3$]
混凝土強度，WGP~建立之校舍耐震能力關係模型參數之一

\item[$P_4$]
校舍長度，WGP~建立之校舍耐震能力關係模型參數之一

\item[$P_5$]
校舍深度，WGP~建立之校舍耐震能力關係模型參數之一

\item[$P_6$]
教室柱數量，WGP~建立之校舍耐震能力關係模型參數之一

\item[$P_7$]
教室柱斷面積，WGP~建立之校舍耐震能力關係模型參數之一

\item[$P_8$]
是否有走廊，WGP~建立之校舍耐震能力關係模型參數之一

\item[$P_9$]
二樓樓地板面積，WGP~建立之校舍耐震能力關係模型參數之一

\item[$P_{10}$]
總樓地板面積，WGP~建立之校舍耐震能力關係模型參數之一

\item[$P_{11}$]
$S_{DS}$，校舍的震區設計水平譜加速度係數，WGP~建立之校舍耐震能力關係模型參數之一

\item[$P_{12}$]
$S_{D1}$，一秒週期的設計譜加速度，WGP~建立之校舍耐震能力關係模型參數之一

\item[$P_{13}$]
教室數量，WGP~建立之校舍耐震能力關係模型參數之一

\item[$P_{14}$]
校舍建築物的跨距，WGP~建立之校舍耐震能力關係模型參數之一

\item[$R^2$]
決定係數

\item[$S_i$]
第~$i$~個輸入參數的重要度

\item[$S_{aD}$]
一般工址或近斷層區域之工址設計水平譜加速度係數

\item[$T$]
建築物基本振動週期

\item[$T_{adj}$]
$C_E$~值修正因子

\item[$T_0^D$]
$S_{D1}/S_{DS}$

\item[$T_{AC}$]
柱等效強度

\item[$T_{AW}$]
牆等效強度

\item[$W_{ij}$]
類神經網路中第~$i$~個輸入層單元與第~$j$~個輸出層單元間之連結權重

\item[$X_i$]
所建構模型中第~$i$~個輸入屬性，也可做為類神經網路中，輸入層第~$i$~個處理單元之輸入值

\item[$Y$]
所建構模型中之輸出屬性

\item[$Y_j$]
類神經網路中輸出層第~$j$~個處理單元之推論輸出值

\item[$k$]
$k$~群交叉驗證初始時，資料的分群數量

\item[$k'$]
Random subsampling~驗證方法中，隨機取樣並建立模型的次數

\item[$net_j$]
類神經網路中輸出層第~$j$~個處理單元之集成函數

\item[$q_1$]
平面及立面對稱性，初步評估調整因子之一

\item[$q_2$]
軟弱層顯著性，初步評估調整因子之一

\item[$q_3$]
裂縫鏽蝕滲水等程度，初步評估調整因子之一

\item[$q_4$]
變形程度，初步評估調整因子之一

\item[$q_5$]
平面耐震性，初步評估調整因子之一

\item[$q_6$]
短柱嚴重性，初步評估調整因子之一

\item[$w_i$]
WGP~決策樹中，第~$i$~條分支線上之權重

\item[$x$]
輸入屬性

\item[$y$]
輸出屬性預測值

\item[$y_i$]
第~$i$~個輸出屬性預測值

\item[$\hat{y}$]
輸出屬性實際值

\item[$\hat{y_i}$]
第~$i$~個輸出屬性實際值


\item[$\alpha$]
命中率的容許誤差

\item[$\beta$]
廣義線性模型中之斜率，也就是迴歸目標

\item[$\epsilon$]
SVM~訓練模型之設定參數

\item[$\theta_j$]
類神經網路中輸出層第~$j$~個處理單元之閥值



\end{SymEntry}


\newpage
\setcounter{page}{1}
\pagenumbering{arabic}
%% 論文本體頁碼回復為阿拉伯數字計頁，並從頭起算
%\pagenumbering{arabic}
%%%%%%%%%%%%%%%%%%%%%%%%%%%%%%%%